\scanpage{209}
\mspage{1}
% archive sheet number: 208

\uline{Catégories multigaloisiennes.}

$\mathcal U$ un univers

$C$ une catégorie

A) \uline{Conditions équivalentes}

(i) Il existe un $\mathcal U$-topos $\mathscr T$, et une équivalence

\[
C\approxeq\mathscr T_{\deleted{\unsure{\mathrm{lcf}}}\unsure{\mathrm{revf}|\mathrm{rev}\,f|\mathrm{cf}}}
\]

où $\mathscr T_{\mathrm{revf}}$ est la \unsure{ss-catégorie|catégorie} pleine de $\mathscr T$ formée des objets loc\textsuperscript{t} ``constants finis''.

(ii) \inserted{$C$ est une $\mathcal U$-catégorie, l'ens des objets mod isom \inserted{de $C$} est $\mathcal U$-petit, et} Il existe une topologie sur $C$, [moins fine que la canonique,] et telle que le foncteur can $C\longrightarrow\widetilde C$ induise une équivalence $C\longrightarrow\widetilde C_{\mathrm{revf}}$.

(ii bis) Itou, en prenant la ``topologie de la descente effective universelle''.

(iii) \inserted{$C$ est une $\mathcal U$-catégorie, dont l'ens des objets mod iso. est $\mathcal U$-petit ;} Pour tout objet $X$ de $C$, \deleted{l'\unsure{ensemble|ens. des|ensemble des}} objets $S$ tels que $X_S$ sur $S$ soit constant fini sur $S$ forment une famille \deleted{couvrante} de descente effective universelle.

(iv) Cf conditions sur feuille 1 bis

B) \uline{Conditions équivalentes}

\deleted{(i) Comme (i) dessus, mais en exigeant de plus} \deleted{que l'objet final de $\mathscr T$ est somme \uline{finie} d'objets} \deleted{connexes}

\deleted{(ii) (ii bis) (iii) Comme plus haut, mais en} \deleted{exigeant que l'objet final de $C$ (ou encore,} \deleted{tout objet de $C$, au choix) est somme finie} \deleted{d'objets connexes.}

\deleted{(iv) $C$ satisfait aux conditions suivantes.}

(i) $C$ satisfait aux conditions de A), et son objet final est somme effective universelle d'objets connexes.

